\section{Semantic feedback and hierarchical credit}

A scalar reward says whether a trajectory worked but discards the explanation. A PPO
critic estimates scalar values at token states; its derived advantage says whether a
sampled action was better or worse than expected, but does not name the mistake or propose
an alternative. Textual feedback can carry that missing structure: which assumption
failed, which tool was available, or which subgoal should change. The human-learning
analogy is instruction rather than a score: one grounded explanation can change several
downstream choices. OPSD converts this semantic intervention into a token-distribution
target rather than compressing it immediately to one number.

Long-horizon credit assignment should therefore be hierarchical. \emph{Trajectory-level}
verification asks whether the complete task succeeded. \emph{Trace-level} diagnosis finds
the responsible model call, subgoal, tool interaction, or context segment.
\emph{Token-level} optimization changes the action distribution inside that selected
trace. These levels are complementary: terminal evidence anchors correctness, a grounded
hint identifies the causal correction, and OPSD or SFT turns that correction into a
gradient. Applying token-level feedback everywhere without first locating the responsible
trace can be as noisy as broadcasting one terminal reward over the whole trajectory.

This suggests an automated post-training loop: verify the outcome, diagnose a trace,
generate and validate semantic feedback, then either convert it into a process reward,
distill its hint-conditioned distribution, or train an explicit
hint-and-correction sequence. Resampling tests whether the behavior actually changed. The
central claim is not that text avoids numerical optimization, but that it preserves a
much richer description until the last step where it must become a gradient. Repeating
this loop is a plausible bridge from post-training to continual agent improvement.
